\def\thetagen{\theta_{\mathrm{True}}}
\def\xbar{\tilde{\x}}

In this section, we illustrate verification of our assumptions and the
consistency of the IJ covariance with a simple closed form example, that of
inferring a univariate normal mean with a misspecified variance. More elaborate
examples can be found in \appref{simple_examples_appendix}, including a Poisson
model in the presence of overdispersed data and multilinear regression in the
presence of heteroskedasticity.

For some fixed $\thetagen \in \rdom{}$ and $\sigma^2 \ne 1$, let us
consider the following data generating process and model:
%
\begin{align*}
%
\textrm{Data generating process:}\quad&&
x_n &\iid \normdist(\cdot | \thetagen, \sigma^2) \\
\textrm{Model:}\quad&&
x_n \vert \theta &\iid \normdist(\cdot | \theta, 1) &
\theta &\sim \normdist(\cdot | 0, 1).
%
\end{align*}
%
We will take $\theta$ to be our quantity of interest, so $\g(\theta) = \theta$.
Since $\sigma^2 \ne 1$ the model is misspecified, and so we expect the posterior
variance $\cov{\post}{\theta}$ to differ from $\gcovtrue$, even asymptotically.

For this simple model, the posterior is available in closed form:
%
\begin{align*}
    \post = \normdist\left( \theta \vert \xbar, (N + 1)^{-1} \right)
    \quad\textrm{where}\quad
    \xbar :={} \frac{1}{N+1} \sumn \x_n.
\end{align*}
%
From this, we can see that $\thetatrue = \lim_{N \rightarrow \infty}
\expect{\post}{\theta} = \thetagen$, It happens in this simple example that
$\thetatrue$ is ``consistent'' for the true mean $\expect{\fdist(\xn)}{\xn}$,
though in general this need not be the case.  For example, the model's trivial
``estimate'' of $\var{\fdist(\xn)}{\xn}$ is $1$, which is never equal to the
true value of $\sigma^2$.

One can readily verify that \assuref{bayes_clt,bclt_okay} are satisfied for this
model and this particular value of $\thetatrue$. Arguably the most challenging
condition is the final part of \assuref{bayes_clt} \itemref{strict_opt}, which
one can show via strong convexity of the log likelihood, together with
convergence of $\frac{1}{N} \sumn \ellgrad{1}(\x_n \vert \thetatrue) \plim 0$ by
the ordinary law of large numbers.  Using special structure of a particular
problem (like strong convexity) to establish a condition like
\assuref{bayes_clt} \itemref{strict_opt} is common in the literature --- see,
for example, textbook proofs of the consistency of M-estimators on unbounded
domains (e.g.~Theorem 9.11 of \citet{keener:2010:theoretical} or Chapter 6.2 of
\citep{huber:2011:robust}).

We begin by establishing \corref{bayes_expectation_dist} and identifying our target,
$\gcovtrue$. By the ordinary CLT applied to our closed-form expression for
$\expect{\post}{\theta}$, we have that
$$
\sqrt{N}\left(\expect{\post}{\theta} - \thetatrue\right) \dlim
\normdist\left(\cdot \vert 0, \sigma^2 \right),
$$
from which we see that $\gcovtrue = \sigma^2$.  It happens, in this case, that
the variance of the limiting distribution is also the limit of the variance,
since
$$
N \var{\xfdist}{\expect{\post}{\theta}} =
    \frac{N^2}{(N+1)^2}\sigma^2 \rightarrow \sigma^2 \gcovtrue.
$$
The asymptotic equivalence between the variance of the limit and the limit of
the variance occurs in this case because $\var{\xfdist}{\expect{\post}{\theta}}$
is uniformly integrable as a sequence in $N$ (Theorem 10.3.6 of
\citet{dudley:2018:realanalysis}).  We emphasize, however, that our theorems do
not assume uniform integrability, and we aim to estimate the variance of the
limiting normal distribution, not the limiting value of the variance, as
discussed in the text surrounding \eqref{gtrue}.

To complete \corref{bayes_expectation_dist}, we need to show that the expression
for $\gcovtrue$ is correct. In this case, $\ggrad{1}(\theta) = 1$, and the
remaining quantities are (up to constants not depending on
$\theta$):
\begin{align}
%
\likhat(\x \vert \theta) ={}&
    -\frac{1}{2N} \sumn \left(x_n^2 - 2x_n \theta + \theta^2\right) + \textrm{Constant}
\quad\textrm{and}\quad
% \nonumber\\
\log \pi(\theta) ={}& -\frac{1}{2}\theta^2 + \textrm{Constant} \Rightarrow \nonumber \\
\thetahat ={}& \xbar
\quad\textrm{,}\quad
\ellgrad{1}(\x \vert \theta) ={} x_n - \theta
\quad\textrm{and}\quad
\likhatk{2}(\x \vert \theta) ={} -1 \Rightarrow \nonumber \\
\scorecovhat ={}& \meann \left(x_n - \xbar\right)^2
\quad\textrm{,}\quad
\infohat ={} 1
\quad\textrm{and}\quad
\gcovhat ={} \infohat^{-1} \scorecovhat \infohat^{-1} = \scorecovhat.
\eqlabel{normal_gcovhat}
%
\end{align}
%
It is readily verified that $\scorecovhat \plim \sigma^2 = \gcovtrue$,
so we have verified the conclusion of \corref{bayes_expectation_dist}.

Finally, we manually verify \thmref{ij_consistent}. The influence score is given by:
%
\begin{align*}
%
\infl_n ={}& N \cov{\post}{
    \theta, -\frac{1}{2} \left(\x_n^2 - 2\x_n \theta + \theta^2 \right)} \\
={}&
 -\frac{1}{2}  N \cov{\post}{\theta, 1} \x_n^2 +
 N \cov{\post}{\theta, \theta} \x_n  + N  \cov{\post}{\theta, \theta^2}
={} \frac{N}{N+1} \x_n + \textrm{Constant},
%
\end{align*}
%
where the constant in the final line is the same for all $\x_n$.  Since the IJ covariance
depends only on the difference $\infl_n - \bar{\infl}$, quantities that are the
same for each $n$ can be neglected in the IJ covariance computation. Therefore
the IJ covariance estimate is given by
%
\begin{align*}
%
\gcovij ={}& \meann \left(\infl_n - \bar{\infl} \right)^2
={} \frac{N^2}{(N + 1)^2} \scorecovhat \plim \sigma^2 = \gcovtrue.
%
\end{align*}
%
Thus $\gcovij$ differs from the $\gcovhat$ given in \eqref{normal_gcovhat} by a
factor that vanishes asymptotically.  So, like $\gcovhat$, $\gcovij$ is a
consistent estimator of the asymptotic variance $\gcovtrue$.
